\exnumber{00}
\extitle{Project 4: A table !}
\exfiles{ex00/zip archive, named Name\_FirstName\_DeliverableName\_Date\_VersionNumber}
\exauthorize{Those named in the exercise}
\makeheaderfiles

Design a Food delivery app (interfaces and user journey for (1) the restaurant owners, (2) the clients, (3) the delivery men (4) the delivery men managers, on mobile and desktop.

DeliverNow is a company that helps restaurants deliver take-away food, thanks to a series of runners.

They recently had bad comments from customers:

'Our customers currently face problems with wrong estimated delivery times, the amount of unnecessary cutlery and condiments.

Our runners complain about the allocation of routes among them.

The managers can't manage their drivers well since they don't have a clear enough vision of the runners.

Finally, the restaurants complain because they have a hard time managing the preparation of the food on-site and take-away.'

Thus, they ask your team to redesign the full experience:

\begin{itemize}
\item For customers
\item For restaurants
\item For runners
\item For runners' managers
\end{itemize}

You have to meet your client after 3 weeks to give a first presentation of your work, after completing a tested mid-fidelity prototype for each user 4 journeys.

The final presentation of your work will take place 2 weeks after the mid presentation, where you'll show a tested high-fidelity prototype of each of the 4 journeys.

\textbf{Estimated time:} 5 weeks

\textbf{Team size:} max 3 people

\textbf{Guidelines:}

\begin{itemize}
\item You may consider building your UX and UI strategy before jumping in. Remember to question the initial brief with user research.
\item Beware all the steps of the design process: Empathy, Define, Problem Statement, Ideation, Prototype, Test.
\item In this context, 4 users = 4 personas = 4 user journeys = 4 prototypes
\item You must interview or ask real people during the Empathy / User Research and Test phases.
\item In your prototypes, if necessary, show your audience the user journey for each persona when there is one order.
\item In both phases, you must test your prototypes with at least 5 people.
\item Make both presentations as if you were showing your work to DeliverNow team. Make sure there are at least 5 people in the audience.
\item Prepare a feedback form to send the audience. The form must include at least the questions listed below.
\item Timing: 15 minutes for the presentation and 5 minutes for Q\&A.
\item You may use the presentation audits to structure your presentations.
\item Upload the documents on Github, in a zip folder named "ProjectTitle\_Name\_FirstName", with all project deliverables as follows: "Name\_FirstName\_DeliverableName\_Date\_VersionNumber".  For example, the first deliverable can be named as "Doe\_John\_MidFiPrototype\_05242024\_V1".
\item Add a title within all written documents.
\item Pay attention to the global aspect of the deliverables. They must be clear, simple and easy to read. You can get inspiration on canva.com but don't overload your design with too much details!
\item Have fun!
\end{itemize}

\textbf{Tools:}

\begin{itemize}
\item Google Forms \url{https://docs.google.com/forms/u/0/}
\item Airtable \url{https://www.airtable.com/}
\item Typeform \url{https://www.typeform.com/}
\end{itemize}

\textbf{Tips:}

\begin{itemize}
\item Effective presentation skills \url{https://academic.oup.com/femsle/article/364/24/fnx235/4587905}
\end{itemize}

Intermediate deliverables:

\textbf{Quote:}

"Your mission in any presentation is to inform, educate, and inspire." Robert Ballard

\medskip

\textbf{Phase 1: Mid-project presentation audit}

Make sure:

\begin{itemize}
\item At least 5 people are attending the presentation
\item The presentation is about the whole UX process, including:
\begin{itemize}
\item The intial brief
\item A project planning, including the UX phases, the tools chosen, the timeline
\item User research findings from qualitative and quantitative research (persona, figures, etc.) for the 4 users
\item 4 problem statements to adapt the initial brief based on the research findings
\item An introduction of each concept that came out of the ideation
\item The animated mid-fidelity prototype (a demo-video) with the 4 user journeys
\item The main feedback from usability testings (5x4)
\item The next steps of the project
\end{itemize}
\item The presentation does not exceed 15 minutes
\item Feedback forms have been shared at the beginning of the presentation
\item You had fun!
\end{itemize}

\medskip

\textbf{Phase 2: Final presentation audit}

Make sure:

\begin{itemize}
\item At least 5 people are attending the presentation
\item The presentation is about the UI process, including:
\begin{itemize}
\item A quick reminder of the previous presentation
\item A project planning, including the UI phases, the tools chosen, the timeline
\item A moodboard
\item A styleguide (colors, branding, typography, some elements of the library)
\item The animated high-fidelity prototype (a demo-video) with the 4 user journeys
\item The main feedback from usability and desirability testing
\item The next steps of the project
\end{itemize}
\item The presentation does not exceed 15 minutes
\item Feedback forms have been shared at the beginning of the presentation
\item You had fun!
\end{itemize}

\textbf{Questions to ask in the feedback forms}

\begin{itemize}
\item What do you remember of this presentation?
\item On a scale from 0 to 10, how clear was the speech?
\item Why?
\item On a scale from 0 to 10, how engaging was the presentation?
\item Why?
\item On a scale from 0 to 10, how coherent was the presentation?
\item Why?
\item On a scale from 0 to 10, how impactful were the visuals?
\item Why?
\item On a scale from 0 to 10, how clear was the problem?
\item Why ?
\item On a scale from 0 to 10, how well does the prototype answers the problem?
\item Why?
\item What recommendation would you like to make to improve this presentation?
\item Why?
\end{itemize}

